\section{选题背景与意义}

\subsection{选题背景}

毛泽东在为《中山县新平乡第九农业生产合作社的青年突击队》所写的按语中指出：“青年是整个社会力量中的一部分最积极最有生气的力量。”\footcite{Mao1977Youth}这句话强调了青年参与社会实践的积极性。大学生的价值取向既关乎个人成长，也关系到未来社会的精神面貌\footcite{Xi2014}。引导大学生扣好人生的“第一粒扣子”、自觉树立和践行社会主义核心价值观，是高校落实立德树人根本任务、培养德智体美劳全面发展的社会主义建设者和接班人的关键所在\footcite{CPCC2017_Ipe,Xi2018_Speech}。近年来，党和国家对贴近青年实际加强价值观教育作出一系列重要部署。2021年，《关于新时代加强和改进思想政治工作的意见》明确要求把思想政治工作贯通于教育教学与日常管理全过程，贴近青年实际做好引导与关怀\footcite{CPCC2021_Yijian}；2024年，党的二十届三中全会对完善立德树人机制、健全培育和践行社会主义核心价值观制度机制作出系统安排\footcite{CPCC2024_Jueding}；2025年印发的《教育强国建设规划纲要（2024—2035年）》进一步提出坚持显性教育与隐性教育相统一，拓展实践育人与网络育人空间\footcite{CPCC2025_Gangyao}。马克思在《关于费尔巴哈的提纲》第六条中指出：“人的本质不是单个人所固有的抽象物，在其现实性上，它是一切社会关系的总和。”\footcite{MarxEngels2009}这一论述为本研究提供了考察视角，即从大学生所处的具体社会关系出发，分析日常交往如何影响其价值判断。大学生正处于价值观确立与定型的关键期\footcite{CPCC2017_Plan}，他们对马克思主义价值观和社会主义核心价值观的理解与认同，也会受到日常人际交往经验的影响。

随着经济社会发展、社会分工细化以及移动互联网普及，当代大学生的日常交往方式正在发生明显变化。在家庭、行政班级和宿舍等长期稳定的熟人圈子之外，越来越多大学生依托线上线下渠道，围绕共同兴趣、具体任务或阶段性需求，与不同专业、不同高校乃至校外的个体建立联系，例如自习搭子、临时项目伙伴、兴趣社群成员、跨校求职互助同伴以及平时联系较少的远亲等。本文所说的弱连接交往，要求双方能够辨认交往对象、有实际回应，并在一定时期内保持联系。只阅读公开帖子、没有与发帖者互动的情况不在此列。这类交往通常不需要较深的情感投入，双方对彼此背景的了解和承担的长期责任也较为有限。客观来看，弱连接交往能够打破固定熟人圈子的局限，帮助大学生获取不同背景的信息与生活经验，在拓宽认知视野、便利日常协作和明确人际边界方面具有一定积极作用。

然而，当这种情感投入较少、关系较易中断的交往方式逐步替代需要长期投入的交往时，可能带来以下问题。具体而言：一方面，在彼此缺乏全面了解的情况下，交往对象集中展示的片面消费片段、升学求职结果或情绪化看法，容易诱发部分学生的攀比焦虑、急功近利倾向以及看待公共事务时的片面极化心态；另一方面，按需求建立联系、任务结束后即停止往来、较少承担长期义务的相处习惯，容易使部分学生把日常人际往来简化为计算得失的实用交换，从而可能削弱室友和同学之间的信任与归属感，也可能使原本联系较少的亲属进一步疏远。此外，若学生将更多时间用于即时性社交消遣，留给深度阅读、学习钻研和自我修养的时间就可能减少。这些倾向在日常点滴中不断累积，与马克思主义关于人的全面发展、崇尚诚实劳动和集体互助的价值立场产生冲突，也与社会主义核心价值观所倡导的爱国、敬业、诚信、友善等育人目标形成现实偏差。列宁在《青年团的任务》中谈到青年学习与共同劳动的关系，提出青年应当安排好学习任务，使自己“每天都能实际完成共同劳动中的某种任务”\footcite{Lenin2017}。这一论述启发本研究关注学生在共同活动中学习合作、履行责任的机会。反观当前高校思想政治工作，日常管理与价值引导仍主要依托院系、班级和宿舍等固定组织展开，面对跨越校园边界、按即时需求聚散的弱连接交往网络，存在感知不够及时、覆盖存在盲区、引导方式针对性不强等问题，难以及时回应学生在流动社交中产生的价值困惑。

在这一背景下，为有效化解弱连接交往带来的负面影响、切实提升新时代高校价值观教育的实效性，本研究立足高校立德树人根本任务，以弱连接交往对大学生价值观的影响为研究对象，系统梳理当代大学生弱连接交往的现实形态与基本特征，深入剖析这种交往方式影响大学生日常价值判断的具体微观过程，揭示其对马克思主义价值观及其当代表现——社会主义核心价值观培育造成的现实冲击，并结合网络文本分析与理论研究探寻切实可行的应对策略，以期为高校在流动社交环境中做好青年价值引领提供支撑。

\subsection{选题意义}

\subsubsection{理论意义}

第一，从日常交往方式变迁切入，系统剖析弱连接交往影响大学生价值观的具体微观过程，有利于深化思想政治教育学对青年价值观形成规律的理论认识。既有学术研究多各自展开：社会学和传播学领域侧重考察弱连接在传递信息、促进求职或提供情绪慰藉方面的实用功能，较少探讨其如何改变青年的价值观念；思想政治教育领域偏重主流价值观的宏观阐释和网络不良思潮的内容剖析，较少深入分析浅层交往本身如何一步步影响大学生的日常价值判断。本研究以马克思在《关于费尔巴哈的提纲》中关于人的本质与社会关系的论述为出发点，把日常交往结构变化同青年价值观念演变紧密联系起来，厘清弱连接交往引发价值偏差的具体传导路径与深层成因，既拓展了弱连接理论在青年价值观研究中的应用边界，也丰富了思想政治教育学关于社会交往与价值观互动关系的学理探讨。

第二，紧扣弱连接交往给马克思主义价值观及社会主义核心价值观培育带来的新挑战，有利于丰富新时代高校价值观引导的理论研究。社会主义核心价值观是马克思主义价值观在当代中国的集中体现，其培育和践行需要不断回应社会生活与青年成长的新变化。随着移动互联网普及和浅层交往在高校日益普遍，以往主要立足固定班集体和传统熟人关系形成的价值观教育理论，难以充分解释流动社交环境下学生观念波动的新特点。本研究结合当代大学生的真实交往实践，清晰界定弱连接交往的范围与特征，系统揭示其在日常生活中与马克思主义价值观、社会主义核心价值观培育要求之间的现实张力，有助于推动社会主义核心价值观培育理论进一步贴近青年的微观生活世界。

\subsubsection{实践意义}

第一，通过梳理弱连接交往中大学生价值观偏移的具体表现与内在原因，有利于帮助高校准确把握青年思想动态，为提升思想政治教育的针对性提供现实依据。习近平总书记强调，做好高校思想政治工作要因事而化、因时而进、因势而新\footcite{Xi2022}，思想政治理论课是落实立德树人根本任务的关键课程\footcite{Xi2020}。当前，大学生在日常消费、学业规划、同伴相处、亲情维系以及公共事务讨论中的许多观念困惑，大量孕育于跨校学习互助、兴趣社群交流和校园搭子往来等弱连接场景之中，而高校传统管理方式对此掌握尚不够充分。本研究结合大学生在网络社区中的真实互动与表达，揭示浅层交往背后隐藏的比较焦虑、功利心态与人际疏离问题，能够帮助高校思政课教师和学生工作队伍精准识别学生关切，把青年在现实交往中遭遇的真问题引入课堂教学与日常引导，切实落实《关于深化新时代学校思想政治理论课改革创新的若干意见》的部署要求\footcite{CPCC2019_Yijian}。

第二，针对弱连接交往的具体作用过程与成因构建系统应对策略，有利于为高校优化育人机制、引导大学生践行社会主义核心价值观提供可操作的实践参考。面对当代大学生普遍重视个人边界、广泛参与弱连接交往的现实，单纯依靠传统的刚性集体管理容易引发逆反，放任浅层社交蔓延又会弱化育人实效。本研究围绕弱连接交往影响价值观的关键环节，从校园实体共同体建设、思政课教学引导、家校协同育人以及校外实践平台拓展等方面提出系统应对路径，探索将主流价值融入青年易于理解与参与的日常生活实践的具体方式\footcite{Shen2021_narrative}，能够为高校党团组织建设、学生社团指导和社会实践开展提供方向参考，引导大学生理性看待浅层社交、克服功利与疏离倾向，在融入集体与服务社会的实践中实现全面发展。

\section{国内外研究现状}
\label{sec:research-status}

围绕“弱连接交往”与“大学生价值观”两个核心主题，本章按国内与国外两条脉络梳理既有研究成果。每一部分均先回顾弱连接交往的研究，重点说明弱关系如何传递信息、数字技术如何改变联系门槛以及当代青年在浅层社交中的真实体验；再回顾大学生价值观的研究，依次归纳青年价值观的现实特点、背后的影响因素以及高校的引导办法。最后综合评述中外研究的进展与不足，明确本研究的切入空间。

\subsection{国内研究现状}
\label{sec:domestic-research}

国内研究扎根于中国从传统熟人社会向现代流动社会转型的现实背景，并随着移动互联网普及和高校思政工作推进不断深入。整体来看，学者既关注弱关系在资源获取与校园“搭子”交往中的新特点，也持续追踪当代大学生价值观的变化趋势及其教育引导路径。

\subsubsection{对弱连接交往的研究}

国内关于弱连接交往的研究，大致沿着从传统熟人社会的关系结构分析，向数字媒介连接方式考察，再向当代大学生“搭子”与轻社交实践延伸的轨迹推进。

\paragraph{强弱关系与资源获取研究}
理解中国社会的弱连接，首先要回到本土人际关系的转型背景。费孝通在《乡土中国》中提出“差序格局”，指出传统乡土社会以血缘和地缘为纽带，人们主要依靠亲戚、邻里等熟人强关系获取生活资源与安全感，圈子内部情感深厚、人情往来密集，但对外较为封闭\footcite{Fei2012}。随着市场化改革与人口跨区域流动加快，封闭的熟人圈子逐渐被打破。边燕杰较早通过对城市居民社会资本的调查发现，在社会交往仍重视人情与信任的背景下，跨越熟人圈子的弱关系开始发挥越来越重要的作用\footcite{ZSHK200403011}；他在后来的关系社会学研究中进一步系统阐明，强关系主要承担信任担保功能，而弱关系擅长在不同群体之间牵线搭桥，能够帮助个体获取熟人圈子之外的新鲜信息与跨阶层职业机会\footcite{Bian2021}。在具体衡量标准上，林聚任在《社会网络分析》中把互动频率、情感深度、亲密程度和互惠内容作为区分关系强弱的四项核心指标，从经验层面证实弱连接的独特价值在于打破小圈子的信息闭塞，成为沟通不同社会群体的桥梁\footcite{Lin2009}。

\paragraph{数字媒介与交往方式研究}
智能手机和算法推荐的普及，改变了人与人建立联系的方式，也成为国内传播学与社会哲学关注的重点。彭兰分析指出，手机带来的“全时在线”状态模糊了工作、社交与私人生活的界限，青年在微信等熟人平台上承受着较重的人情维护和形象管理压力，因而产生收紧朋友圈权限甚至阶段性断联的“反连接”诉求\footcite{GJXW201902002}；她随后进一步提出“液态连接”来概括数字化人际关系的新形态，指出如今的交往不再依赖稳固的熟人外壳，而是跟着具体场景和即时需求快速聚合，事情办完或场景转换便自然散去\footcite{Peng2021}。蓝江从社会哲学角度将这一现象概括为“微粒化社会”，认为数字媒介拆解了家庭、邻里和单位等厚重的实体联系，把具体的人拆散成依靠兴趣标签临时聚合的独立个体，使人际往来逐渐失去共同生活积累下来的约束力\footcite{Lan2022}。胡翼青和王敏则将这一变化归纳为日常交往的“轻量化转向”，指出年轻人在社交中越来越倾向于卸下沉重的人情负担，追求门槛低、维系省力、随时可以无痛退出的互动体验\footcite{Hu2019}。从技术支撑来看，喻国明和曲慧在探讨算法推送机制时指出，智能算法通过对用户日常行为数据的抓取与分析，可以绕过现实熟人关系直接匹配细分兴趣节点，从底层重塑了弱连接的生成途径与信息流动方式\footcite{XJSF202001015}。进一步针对青年网络圈层文化，蔡骐研究发现，靠算法和共同爱好聚起来的兴趣群体往往内部容易产生强烈共鸣，对外却筑起较高的沟通壁垒，反而加大了不同圈层之间相互理解的难度\footcite{Cai2018}。

\paragraph{青年轻社交与搭子交往研究}
线上流动连接延伸到大学校园，催生了饭搭子、自习搭子、运动搭子、拼单搭子等广泛流行的“轻社交”现象。廉思基于全国性青年调查分析指出，“搭子社交”反映出青年从传统班级、社团等固定集体中逐渐抽离后，借助细分领域的简单约定建立起浅层陪伴关系；大学生试图在获得即时陪伴与减少人情责任之间寻找平衡，但长期依赖这种按功能拆分的交往，潜藏着削弱深度沟通意愿、降低共情能力和回避矛盾磨合的隐患\footcite{Lian2024a}。结合升学与就业压力，孙寿涛和张晓芳对Z世代大学生的轻社交行为展开分析，指出这种选择并非单纯冷漠，而是年轻人在学业竞争激烈、未来充满不确定性的背景下，为节省精力、减少人际摩擦和精神内耗而采取的自我保护策略\footcite{Sun2023}。这种心理防御在日常相处中表现为对人际距离的精细把控：刘航分析指出，搭子社交之所以在青年群体中迅速流行，是因为它精准迎合了当代大学生对清晰个人边界、低风险交流和实用互助的需求\footcite{LiuHang2023}；陈姣和牛菲则提醒，如果过度依赖此类浅层关系，容易把原本带有情感温度的同学往来简化为按需调用、即时算清的实用交易\footcite{ChenJiao2024}。邢海燕和邸涵针对大城市青年日常时空边界的调查也印证了这一趋势，表明当代青年习惯主动划清私人生活与公共社交的界线，排斥未经打招呼的熟人介入，借由高度可控的浅层接触来保持个人生活的自主与轻松\footcite{Xing2022}。

\subsubsection{对大学生价值观的研究}

国内学界围绕大学生价值观的研究，紧扣改革开放以来的社会变迁与高校立德树人实践，形成了涵盖现状特点追踪、成因剖析与思政教育引导三个环节的完整脉络。

\paragraph{中国大学生价值观现状研究}
通过长期大样本调查，学者们较清晰地勾勒出中国大学生价值观的演变趋势。黄希庭和郑涌较早立足中国文化与社会现实，从人生目标、职业追求、人际伦理和公共责任等方面构建了适用于中国青年群体的价值观测量体系\footcite{Huang2005}。金盛华和辛志勇在梳理中国人价值观发展趋势时同样指出，改革开放以来，青年价值观逐步从单一集体本位向兼顾个人发展与社会贡献的多元结构演变\footcite{BJSF200303008}。近年来，廉思依托全国大样本调研将当代大学生的价值面貌概括为“宏微双轨”：在国家发展大局、制度认同和民族复兴等宏观立场上态度鲜明且高度稳固，但落到考研、求职、婚恋和日常相处等具体生活选择时，受竞争加剧与回报分化影响，价值取向明显变得务实且分化加剧\footcite{Lian2024b}。沈壮海和王芸婷基于全国大学生思想政治状况滚动调查也发现，受访大学生总体上表现出较强的爱国认同，同时在学业深造、就业压力和人际适应上承受着突出的现实焦虑\footcite{Shen2022}。冯刚和曾耀荣在综述近十年核心价值观培育研究时指出，正是这种宏观认同稳固与微观困惑并存的局面，推动学界从早期的宏观理论宣讲转向如何贴近大学生日常生活、回应具体关切的深入探索\footcite{Feng2022}。在人际信任与亲密关系方面，辛自强和周正通过跨二十年的追踪对比分析发现，伴随人口高频流动与陌生人社会扩大，中国大学生对一般他人的信任度呈明显下降趋势\footcite{XLXD201203004}；辛自强在后续关于青年社会心态的研究中进一步指出，青年在跨圈层交往中的戒备心理和防御边界持续收紧\footcite{Xin2020}。风笑天针对青年婚恋观变迁的专项研究也相互佐证，表明当代大学生愈发看重精神独立、生活节奏自主和边界自由，对需要高度情感投入和长期责任捆绑的传统亲密关系态度更趋谨慎\footcite{FengXiaotian2018}。

\paragraph{大学生价值观影响因素研究}
现有文献主要从社会结构转型和数字媒介环境两条线索，解释大学生价值观发生上述变化的原因。在社会结构层面，阎云翔在《中国社会的个体化》中提出“制度化个体主义”，指出市场化进程和用工制度改革促使青年从大家庭与传统单位的庇护中抽离出来，必须独自面对学业竞争与市场风险，因此个人权利、自主选择与自我负责迅速成为青年做判断时的核心标准\footcite{Yan2012}。从代际成长环境看，周晓虹围绕“文化反哺”与代际转型指出，熟悉互联网的年轻一代在信息获取与文化消费上拥有更强的话语权，改变了传统的长幼单向传授格局\footcite{Zhou2020}；李春玲从代际社会学视角刻画了中国“Z世代”的群体特征，说明成长于物质丰裕与互联网普及期的大学生普遍视野开阔，崇尚平等商量、反感居高临下的说教、尊重个体生活差异\footcite{Li2021}。然而，陆士桢和蔡梦雅在剖析数字时代青年价值观新特征时指出，当青年缺乏稳定的集体支持、同时面临激烈的就业竞争时，不少青年在日常交友、社团活动和择业中更容易滑向看重即时回报与得失算计的实用主义心态\footcite{Lu2022}。在媒介环境层面，胡正荣和李继东指出，移动互联时代去中心化的传播格局削弱了传统权威自上而下传递价值观的影响力\footcite{Hu2017}。隋岩和曹飞从群体传播视角深入分析发现，离开实体班团集体的青年依托网络弱连接形成各类兴趣圈层，圈层内部高频的专属黑话互动与密集情绪感染，会在潜移默化中改写青年的是非标准与价值排序\footcite{Sui2021}。朱哲和刘伟进一步探讨了网络圈层化对大学生核心价值观认同的冲击，认为商业平台算法制造的“信息茧房”助长了消费主义和泛娱乐化风气的扩散\footcite{Zhu2021}；崔聪和张励仁则结合网络民粹思潮分析指出，缺乏事实核查和历史背景支撑的碎片化、情绪化信息流，容易让价值观尚未完全定型的部分学生陷入非黑即白、简单贴标签的偏激判断\footcite{Cui2023}。

\paragraph{大学生价值观的思想政治教育引导研究}
面对个体化趋势与网络浅层社交带来的新挑战，国内思政学界在继承经典原理的基础上，探索了更加精准、贴近日常的育人路径。张耀灿等在《现代思想政治教育学》中阐明，主流价值观转化为学生内心的真实信念，需要经历“认知—认同—内化—外化”的渐进过程，教育内容若能联系学生的交往经验并引起情感共鸣，外在规范就更有可能转化为学生的价值认同\footcite{Zhang2006}。沈壮海从文化视野出发，主张跳出单向说教，通过把宏大叙事转化为贴近青年生命的日常叙事、以润物无声的文化氛围回应学生的精神需求，打通价值认同的现实通道\footcite{Shen2013}。刘建军则围绕社会主义核心价值观“日常化、具体化、形象化、生活化”展开论述，强调要把抽象原则转化为融入课堂、宿舍起居和同伴交往的具体行为准则，引导学生在日常互动中践行诚信与友善\footcite{Liu2016}。在制度层面，教育部等部门先后印发《高校思想政治工作质量提升工程实施纲要》\footcite{MOE2017_Gangyao}、《关于加快构建高校思想政治工作体系的意见》\footcite{MOE2020_Gangshen}以及《全面推进“大思政课”建设的工作方案》\footcite{MOE2022_Fangan}，对“十大育人”体系与社会大课堂融通作出系统部署。顺应这一导向，近年研究重点探索如何走进青年的网络弱连接圈层：骆郁廷和郭莉探讨了数字赋能“精准思政”的内在机理，主张借助大数据和智能技术动态感知学生在线上的社交轨迹与隐性困惑，实现分众化精准供给\footcite{Luo2022}；戴艳军和孙晓晖聚焦话语转化，提出把严肃抽象的“文件话语”“学术话语”转化为贴近青年交流习惯的“生活话语”“互动话语”\footcite{Dai2021}；杨晓慧则围绕数字时代高校网络思政的分众化引导提出，要将线上弱连接圈层的柔性融入与线下校园共同体（如一站式学生社区、科研团队和社会实践）的情感温度重构结合起来，在虚实协同中重新培植大学生的集体归属感与社会责任感\footcite{Yang2023}。

\subsection{国外研究现状}
\label{sec:international-research}

国外对弱连接交往与青年价值观的研究起步较早，横跨社会学、心理学、传播学、计算社会科学与教育学等多个领域，既积累了经典的理论概念，也开展了大量基于真实数据和课堂实验的实证检验。

\subsubsection{对弱连接交往的研究}

国外关于弱连接交往的研究，重点围绕弱关系如何带来信息与机会、互联网如何降低连接成本，以及浅层交往给现代人带来怎样的情绪体验与关系困境三个问题展开。

\paragraph{弱连接与机会获取}
西方古典社会学很早就注意到现代人际关系由紧密走向松散的趋势。滕尼斯（Tönnies）在《共同体与社会》中对比了基于血缘、邻里默契的传统“共同体”和基于理性算计、明确契约的现代“社会”，指出现代化使人际纽带从全面照顾转向特定功能合作\footcite{Tonnies2019}。齐美尔（Simmel）在分析现代大都市生活时指出，分工细化让个体同时处于多个松散交叠的社交圈子里，这既带来了前所未有的行动自由，也拉开了人与人之间的心理距离，催生出都市人理智算计与冷漠防御的心态\footcite{Simmel2002}。真正奠定该领域理论基础的是美国社会学家格兰诺维特（Granovetter），他在经典论文《弱关系的力量》中从互动频率、情感深度、亲密程度和互惠交换四个方面衡量关系强弱，证明联系紧密的强关系大多圈在背景相似的熟人内部，信息高度重复；而联系较少、感情较浅的弱关系能跨越小圈子充当“信息桥”，为人们带来跨行业的求职信息与新思路\footcite{Granovetter1973}。顺着这一思路，伯特（Burt）提出“结构洞”理论，说明处于两个互不相连群体之间的中间人，能够优先获取不同圈子的新信息并掌握协调资源的主动权\footcite{Burt1992}。帕特南（Putnam）则把社会资本区分为凝聚熟人内部的“粘合型社会资本”和连接不同背景人群的“桥接型社会资本”，指出后者对于打破群体偏见、拓宽公共视野和培育广泛社会信任具有关键作用\footcite{Putnam2000}。近五年，计算社会科学用超大规模数据检验了这些经典判断：拉杰库马尔等（Rajkumar et al.）在职业社交平台上对2000万名用户开展长达五年的随机对照实验，证实中等强度的弱连接最能促进青年跨行业求职与职业流动\footcite{Rajkumar2022}；切蒂等（Chetty et al.）追踪分析7200万名社交媒体用户（涵盖全美绝大多数高校学生）后发现，跨越不同收入阶层的弱连接（即“经济连通性”），比小群体内部团结或参加社团更能提升青年的学业抱负与向上流动机会\footcite{Chetty2022}。

\paragraph{数字媒介与弱连接的形成}
互联网与社交媒体大幅压低了建立和维护弱连接的门槛。海桑思韦特（Haythornthwaite）提出“潜在关系”概念，指出大学校园或大型机构里原本就存在大量因同校、同城或同群而具备联系可能、却尚未实际往来的人，数字媒介极大降低了寻找和破冰的成本，让大学生可以按需把这些沉睡的联系随时变成现实弱连接\footcite{Haythornthwaite2005}。埃里森等（Ellison et al.）对美国公立大学本科生展开追踪调查，发现大学生在社交媒体上维系广泛的校园“泛泛之交”，能有效积累桥接型社会资本，尤其有助于性格内向或有社交焦虑的新生适应异地求学生活、增强校园归属感\footcite{Ellison2007}。雷尼和威尔曼（Rainie \& Wellman）将这种数字时代的社交形态概括为“网络化个体主义”，说明年轻人正在从地域和家族等固定捆绑的社群中走出来，转而以自己为中心，根据不同需要灵活编织跨越空间的松散联系网\footcite{Rainie2012}。拜耳等（Bayer et al.）的系统综述进一步揭示了支撑这种网络的日常机制，指出社交平台提供的各类功能设计——比如阅后即焚、分组可见、兴趣群聊和轻量点赞——让年轻人能够精细地控制社交距离、在线状态和自我暴露的深浅\footcite{Bayer2020}。

\paragraph{弱连接的交往体验与关系维系}
关于弱连接给现代人带来的心理体验与长远影响，国外行为心理学和社会批判理论给出了截然不同却彼此互补的观察。在积极一面，桑德斯特伦和邓恩（Sandstrom \& Dunn）通过校园追踪实验发现，大学生在日常生活中与不太熟悉的同学、社团外围成员或咖啡店员进行简短友善的闲聊，就能显著提升当天的愉悦感和归属感\footcite{Sandstrom2014}；阿蒂尔等（Atir et al.）针对陌生人交谈开展系列行为实验进一步证明，人们普遍低估了与陌生人或弱关系对象聊天所能带来的信息收获与好心情，只要迈出搭话的第一步，往往能打破固有局限、获得意外启发\footcite{doi:10.1073/pnas.2206992119}。在隐忧一面，社会理论家深入剖析了浅层连接过度蔓延后对深度关系的侵蚀。吉登斯（Giddens）指出，现代社会让人际关系脱离了传统宗族和地域的外部约束，演化成全靠双方当下心理满意来维系、随时可以单方面退出的“纯粹关系”\footcite{Giddens1998}。鲍曼（Bauman）在《液态之爱》中更为尖锐地批判，消费主义逻辑已经渗入日常交友，诱导现代人像挑选快消品一样建立随时可开启、随时可无痛终止的“轻盈连接”，让人一边贪恋低成本的陪伴，一边本能地逃避需要磨合与承担责任的深度关系\footcite{Bauman2003}。罗萨（Rosa）从时间节奏的角度提出“社会加速”理论，认为科技、社会变化与生活步调不断提速，迫使年轻人在碎片时间里塞满大量浅层社交信息，挤干了静心倾听和建立深层情感共鸣所需的时间，使高频联系的背后弥漫着空虚感\footcite{Rosa2015}。这一判断得到了实证研究的有力印证：特克尔（Turkle）通过长期访谈北美高校学生发现，大学生越来越习惯用可以反复修改、零情绪风险的文字消息替代面对面交流，长此以往导致察言观色的共情能力、深度倾听耐心以及包容不同意见的雅量明显退化\footcite{Turkle2015}；霍尔（Hall）的实证测算更直观地揭示了时间挤占问题——从泛泛之交变成普通朋友平均需要约50小时的真实相处，变成亲密挚友则需要200小时以上的高质量陪伴，大学生把大量闲暇耗在刷动态和点赞等浅层连接上，客观上挤占了培养真挚友谊的稀缺时间\footcite{Hall2019}。

\subsubsection{对大学生及青年价值观的研究}

国外学者结合青年成长阶段与数字时代特征，围绕大学生价值观的测量与特点、环境与媒介成因以及学校教育的引导机制，开展了大量跨文化调查与课堂实验研究。

\paragraph{大学生及青年价值观现状与特点研究}
西方发展心理学和跨文化研究系统刻画了当代大学生的价值特征。阿内特（Arnett）提出“初显期成人”（18—25岁）理论，指出高等教育延长使大学生脱离了未成年时的家庭依附，又尚未步入稳定婚姻和终身职业，处于探索自我、流动性高且高度关注自身发展的独特时期，因此格外看重自主决策和多元体验\footcite{Arnett2000}；他在后续修订著作中进一步证实，这一阶段强调情感自主与个性探索的价值倾向在当代各国青年中愈发普遍\footcite{Arnett2023}。在价值观分类测量方面，早期心理学家斯普朗格（Spranger）从人格整体出发，将人类主导价值划分为理论型、经济型、审美型、社会型、政治型和宗教型六类\footcite{Spranger2017}。在此基础上，施瓦茨（Schwartz）构建了极具国际影响力的“人类基本价值观环形模型”，用“关爱他人还是追求个人成功”“乐于求变还是看重传统稳定”两组对照维度划分出十类基本价值观\footcite{Schwartz1992}；施瓦茨等（Schwartz et al.）在2012年进一步将模型细化为十九类，并基于跨文化样本证实，大学阶段青年在自主思考、乐于尝试和普世包容上的权重显著高于中老年群体，而在遵守传统与顺从规范上得分偏低\footcite{Schwartz2012}。政治社会学家英格尔哈特（Inglehart）的长期追踪也印证了这一趋势，认为成长于物质相对充裕年代的年轻一代，价值重心正从关注温饱与安全转向强调自我表达和生活质量的“后物质主义”\footcite{Inglehart2018}。然而，2010年代智能手机和社交媒体普及之后，情况出现了新的复杂变化。特温格（Twenge）基于全美数百万大中学生跨世代追踪数据发现，当代“手机世代”呈现出“高包容与高脆弱”并存的特征：他们在平等、多元文化等议题上态度开放，但线下同伴聚会、打工实践和独立出行的比例大幅下滑，孤独感与抑郁、焦虑等心理脆弱指标显著上升\footcite{Twenge2017}。卢基安诺夫和海特（Lukianoff \& Haidt）指出，这促使不少大学生把“避免心理不适”奉为高于一切的准则（情感安全主义），习惯把日常观点分歧放大为严重伤害，面对不同意见时倾向于回避或拉黑\footcite{Lukianoff2018}；海特（Haidt）在2024年的最新著作中进一步追溯其根源，论证青年成长方式从线下自由玩耍、在同伴小摩擦中学会协商，全面转向手机屏幕与虚拟弱连接，剥夺了年轻人在真实交往中锻炼抗挫折能力的机会\footcite{Haidt2024}。

\paragraph{价值观影响因素研究}
国外学界从成长环境、代际经历、社会交往规则以及平台微观机制四个层次剖析青年价值观的成因。罗克奇（Rokeach）较早将价值观区分为追求理想人生状态的“终极价值观”和指导日常做法的“工具价值观”，说明社会制度与文化规范主要通过调整这两类价值的轻重排序来影响个人判断\footcite{Rokeach1973}。在环境层级上，布朗芬布伦纳（Bronfenbrenner）提出“人类发展生态学模型”，强调青年价值观在家庭与同伴圈、学校、社区以及宏观社会文化的层层互动中形成\footcite{Bronfenbrenner1979}；顺应移动互联时代，纳瓦罗和图吉（Navarro \& Tudge）提出“新生态理论”作出修正，将青年日常所处的微观环境拆分为面对面的现实空间和跨时空的网络空间，指出社交媒体上的高频弱连接互动已经演化为独立塑造大学生价值认知与自我认同的成长空间\footcite{Navarro2023}。在时间维度上，曼海姆（Mannheim）的代际理论指出，17至25岁是青年脱离原生家庭习惯、初次独立接触社会并形成稳定政治与社会价值观的“关键敏感期”，这一时期共同经历的技术变革与社会事件可能成为影响同代人价值取向的共同经历\footcite{Mannheim1952}。在深层交往规则层面，哈贝马斯（Habermas）在《交往行为理论》中批判了市场和官僚体制中的效率算计不断侵入日常生活的现象，指出这种扩张正在蚕食以真诚沟通和相互理解为核心的正常交往，为审视当代青年社交中的功利算计提供了深刻视角\footcite{Habermas2004}。在此基础上，近年的实证研究查明了网络弱连接扭曲青年价值判断的三重微观机制：其一，桑斯坦（Sunstein）指出，算法个性化推送筑起的“回音室”不断削弱社会成员共享的公共经验基础\footcite{Sunstein2017}，而贝尔等（Bail et al.）通过现场随机实验发现，在参与者之间缺少既有信任的社交媒体网络中，单靠算法让用户接触对立阵营的陌生人观点，反而容易激发强烈的群体防御心理，加剧立场极端化\footcite{Bail2018}；其二，布雷迪等（Brady et al.）对海量社交推文的大数据分析表明，在松散的弱连接传播网络中，带有激烈道德指责与愤怒情绪的言论比温和说理更容易获得点赞和转发，平台的可见性奖励不断诱导青年在公共议题上抛弃全面考量，采取更具攻击性的偏激立场\footcite{Brady2021}；其三，韦尔杜因等（Verduyn et al.）的系统综述证实，由于线上弱连接网络充斥着泛泛之交精心修饰的“高光生活”，大学生在缺乏真实背景了解的情况下被动浏览这些动态，极易陷入无休止的向上社会比较、心理落差与“错失恐惧”，潜移默化地助长物质攀比并削弱自我价值感\footcite{Verduyn2020}。

\paragraph{教育对价值观的塑造作用研究}
针对个体化与数字化带来的青年价值失范风险，西方教育学、道德心理学和媒介素养学者探索了多种具体干预路径。在集体生活与实践育人方面，涂尔干（Durkheim）在《道德教育》中指出，学校道德教育的关键在于培养纪律精神、关心集体的利他情怀和审慎理解规则的自主能力，强调有组织的学校集体生活有助于减轻个人与群体疏离所带来的利己倾向\footcite{Durkheim2001}；杜威（Dewey）则提出“学校即社会”，批评脱离实际交往的抽象说教，论证真正有生命力的公共责任感只能生长于学生共同合作解决实际问题的日常交往之中\footcite{Dewey2013}。在道德心理与情感培育方面，科尔伯格（Kohlberg）通过长期追踪提出道德认知发展的六阶段理论，证明引导青年超越功利交换迈向公正原则的关键，在于教师有意识地组织“道德两难讨论”，借助同伴之间的真实观点碰撞打破自我中心思维\footcite{Kohlberg1981}；诺丁斯（Noddings）从“关怀伦理学”角度补充指出，化解现代人际冷漠不能只靠抽象推理，还要把教育重心落回具体的人际关怀，通过教师的榜样示范、真诚平等的对话、制度化的关怀实践以及对学生善意动机的肯定，培养青年体谅他人的共情能力与责任感\footcite{Noddings2013}。在高校实证评估方面，梅休等（Mayhew et al.）整合了21世纪以来全美数千项高校追踪调查，证实大学之所以能显著提升学生的利他精神、公共责任感与包容度，最核心的推动力来自课堂之外高参与度的同伴真实合作、跨群体项目、住宿学院共同体生活以及深入社区的“服务性学习”，这些深度协作实践能够有效对冲浅层连接带来的疏离与功利心态\footcite{Mayhew2016}。面对网络弱连接环境中的算法偏见与信息误导，崔等（Choi et al.）的调查表明，教师自身的网络实践经验与数字伦理意识，关系到学校数字公民课程能否回应学生在网络交往中遇到的问题\footcite{Choi2018}；温伯格等（Wineburg et al.）近年在多所高校开展的课堂对照实验进一步证明，在通识课中系统教授专业事实核查员使用的“横向阅读”方法——即遇到陌生网络信息时主动跳出当前页面去核查信源背景与利益关联——能够显著提高大学生辨别误导信息的能力，培养审慎求真、不盲从的理性判断习惯\footcite{Wineburg2022}。

\subsection{国内外研究综合述评}
\label{sec:research-review}

综观国内外既有研究，学界围绕弱连接交往演变与大学生价值观变迁已取得丰硕成果，为本研究奠定了扎实的理论与方法基础。已有多项国内外研究指出，在社会个体化转型与移动互联网普及的共同推动下，以情感投入浅、边界清晰、退出成本低和按需匹配为特征的弱连接交往（包括校园“搭子”、轻社交以及线上跨圈层联系），已成为当代大学生日常互动的重要方式。既有文献通过大规模数据实验、长期追踪调查和深度访谈，既讨论了弱连接在拓宽信息来源、助力职业流动、缓解新生适应焦虑和提供轻松情绪慰藉方面的积极作用，也指出了浅层连接过度蔓延背后潜藏的人际功利化、社会信任下降、共情能力退化、向上比较焦虑以及公共讨论极化等隐忧。

总体而言，国外研究长于用大规模数据和课堂对照实验检验因果机制与干预效果；国内研究则紧密立足中国从熟人社会向流动社会转型的本土脉络，精准把握当代大学生价值观“宏微双轨”的现实特点，并依托“十大育人”体系、“精准思政”和“大思政课”探索价值引领路径。

然而，结合本研究的主题加以审视，既有文献在以下三个方面仍有进一步深化的空间。第一，在理论视角上，社会网络分析与思想政治教育价值引导仍相对割裂。社会学和传播学关于弱连接的研究多聚焦求职信息传递、社会资本积累或即时情绪获取的实用功能，往往把行动者的价值观念当成既定前提，抑或停留在宏观现代性批判层面；思想政治教育领域的网络研究则多集中于网络不良思潮辨析和话语转换技巧，较少从马克思在《关于费尔巴哈的提纲》中提出的人的本质与社会关系的观点出发，系统分析弱连接交往中的相处规则如何影响大学生的日常价值判断，并由此给高校价值观教育带来哪些挑战。

第二，在作用过程上，缺少跨生活领域的整体整合。受学科分工限制，现有实证研究大多只截取青年生活的某一个切面（例如单独讨论“搭子社交”的心理防御，或单独检验社交媒体的向上比较），尚未把大学生的个人消费休闲、学业职业追求、同伴及亲属往来与社会公共参与作为一个相互牵动的整体加以考察，也未能充分揭示信息片面展示、低义务相处习惯、浅层满足替代与松散群体环境是如何共同作用于日常价值判断的。

第三，在应对策略上，针对流动社交特点的系统方案尚显不足。既有引导措施容易走向两个极端：要么依赖传统行政班级与刚性集体活动去约束学生的浅层社交，因忽视青年合理的边界诉求而引发逆反；要么一味顺应算法推荐去迎合碎片化、轻量化偏好，无形中强化了工具算计与圈层封闭。学界仍缺乏一套既尊重当代青年交往特点、又能通过重构校园实体协作共同体和强化马克思主义方法论训练来化解价值偏差的系统应对路径。这正是本研究力图推进和突破的方向。

\section{学位论文研究方案}

\subsection{研究目标}

围绕弱连接交往对大学生价值观的影响及其应对策略这一主题，本研究聚焦三项关键问题并预期达到以下目标：

第一，明确弱连接交往的判定标准并构建大学生价值观四维评判标准。为区分具有一定延续性的人际联系与围绕公共话题发生的一次性互动，并将抽象价值规范转化为可观察的指标，本研究将研读马克思主义经典著作、党和国家重要政策文献以及国内外相关学术成果，明确当代大学生弱连接交往的范围及其基本特征，结合社会主义核心价值观的育人要求与公民道德生活领域的划分，构建考察大学生价值观变化的四维分析框架与具体评判标尺。

第二，揭示弱连接交往引发大学生价值观负面偏移的微观作用过程与现实成因。针对浅层社交规则如何一步步转化为青年日常价值判断变化尚欠清晰的问题，通过收集与分析多平台大学生网络文本资料，梳理弱连接交往背景下大学生在个人生活、人际交往、学习劳动与择业、社会公共四个方面的价值观负面偏移表现，打通从日常浅层社交规则到具体价值判断变化的传导链条，剖析其背后的深层现实成因，揭示其对马克思主义价值观及社会主义核心价值观培育造成的现实冲击。

第三，破解流动社交环境下高校思想政治教育感知滞后与引导脱节的难题。针对传统固定班团管理覆盖存在盲区、难以兼顾学生合理个人边界与主流价值引领的问题，立足高校思想政治教育主责，围绕弱连接交往的特点与关键影响环节，从校园实体共同体建设、弱连接资源反向利用、思政课教学转化和校社协同联动等方面提出切实可行的引导方案，为高校在流动社交环境中提升价值观教育实效、落实立德树人根本任务提供参考。

\subsection{研究内容}

本研究立足新时代高校立德树人根本任务，以弱连接交往对大学生价值观的影响为核心对象，围绕理论框架构建、现实表现考察、过程成因剖析与应对策略提出四个层面依次推进。

首先进行理论梳理与分析框架构建。以马克思在《关于费尔巴哈的提纲》中关于人的本质与社会关系的论述为理论起点，并参考《思想政治教育学原理》对相关问题的阐释\footcite{MarxEngels2009,sizheng2018}，结合社会网络理论，以“对象可辨认、有实质回应、有一定延续性”界定弱连接交往边界，涵盖兴趣社群、跨校备考求职互助、校园“搭子”及平时联系较少的亲属往来，排除只浏览而未回应的情况，以及参与者之间未形成持续联系的一次性争论。同时，承接《新时代公民道德建设实施纲要》的社会生活分域\footcite{gangyao2019}，参照马工程教材《思想道德与法治》\footcite{sixiangdaode2023}、《伦理学》\footcite{lunlixue2020}及习近平总书记“核心价值观，其实就是一种德”的论断\footcite{xi2017wenhua}，结合在校大学生成长实际，确立个人生活、人际交往、学习劳动与择业、社会公共四维价值观分析框架，对照马克思主义价值观与社会主义核心价值观制定评判标尺。

其次考察弱连接交往中大学生价值观负面偏移的现实表现。依托多平台大学生网络文本资料，还原在校生在浅层社交中的真实心态与行为倾向：既捕捉攀比消费、表层消遣以及将专业学习窄化为绩点与文凭筹码的功利浮躁，也聚焦“搭子”式边界防御、宿舍集体归属感流失与刻意疏离远亲等关系淡化现象，同时关注松散社群中以情绪化表态替代讲理协商、公共规则意识与公益热情弱化的问题，呈现这些日常偏差与社会主义核心价值观育人要求之间的现实差距。

随后剖析微观作用过程与深层成因。结合网络文本加以印证，重点拆解浅层互动规则改写价值判断的四条内在路径：片面展示光鲜消费与竞争结果诱发的向上比较焦虑；受市场等价交换观念影响，按需匹配、随时散场的低义务相处习惯向同伴与亲属关系延伸；省力即时的浅层满足对深度阅读、真挚友谊与亲情沟通的时间挤占；以及缺乏稳定共同劳动的松散群体环境助长的情绪放大与公共责任淡化。在此基础上，从社会分工细化、升学就业竞争压力、家庭空间分离与高校传统班团管理覆盖局限探寻现实根源，并结合现实强关系支持差异说明负面影响的触发条件与适用限度。

最后提出高校思想政治教育应对策略。紧扣前述过程与成因，在尊重青年合理个人边界的前提下强化主流价值引领：依托党团支部、科研团队与实践项目做实校园实体集体，借助家校协同与返乡实践缓和亲属疏离；反向利用弱连接跨界优势引入行业一线奋斗者，并通过读书共同体与志愿服务实施分类纠偏；强化思政课中马克思主义社会关系、劳动与消费观的方法论训练，推动核心价值观融入日常叙事；联动校外社区实践基地与校友导师网络，构建校社协同育人格局。

% 【此处保留原有的“毕业论文框架目录”部分】

\subsection{拟采用的研究方法与技术路线}

本研究坚持辩证唯物主义和历史唯物主义的根本方法，将弱连接交往对大学生价值观的影响置于经济社会变迁与校园现实竞争压力的物质条件中加以考察，既客观承认其在拓宽信息渠道方面的积极作用，也重点剖析其浅层互动规则诱发观念偏差的消极影响。围绕“提出问题—分析问题—解决问题”的基本逻辑，研究将文献理论分析、网络文本挖掘与质性内容分析有机衔接起来，先界定研究对象，再分析网络材料所呈现的交往过程，最后据此提出育人对策。

研究首先从文献研究与理论分析入手，为全篇奠定学理根基。通过系统精读《马克思恩格斯文集》中的《1844年经济学哲学手稿》《关于费尔巴哈的提纲》《德意志意识形态》《共产党宣言》《资本论》（第一卷）等经典文本，梳理马克思主义关于人的社会关系本质、交往方式反作用于思想观念、劳动观、消费观及人的全面发展的核心论述；同时研读党的创新理论文献、高校思想政治工作政策文件及马工程重点教材，明确考察大学生四维价值观变化的价值准绳。在此基础上，广泛检索国内外关于弱关系理论、轻社交、搭子交往、青年亲属疏离及大学生价值观演变的代表性研究，立足马克思主义立场对社会学、传播学与思想政治教育学成果加以整合，划定弱连接交往的判定边界，确立个人生活、人际交往、学习劳动与择业、社会公共四维分析框架与负面偏移评判标准。

在确立分析框架之后，研究转向经验考察，运用网络文本挖掘法从自然交往场景中捕捉大学生关于弱连接互动的真实表达，并从宏观层面把握高频话题与话语关联。为避免开放网络社区中用户身份混杂、泛泛而谈的问题，数据采集采取平台分层、高校身份核验与弱连接情境限定相结合的筛选策略：以代表性高校公开论坛作为身份校准样本，重点覆盖大学生活跃的小红书、豆瓣高校及青年小组、知乎高等教育及应届生板块近三年的公开主题帖与互动评论；通过高校专属标签、校内板块属性及发帖者自述的本科生、研究生、室友、保研、春招等在校信息过滤非大学生群体，严格按照“对象可辨认、有实质回应、有一定延续性”的标准，定向收集校园拼单与搭子互动复盘、轻社交边界讨论、减少亲戚往来反思、考研与春招打卡互助（围绕“绩点”“内卷”“上岸”等竞争议题展开）以及校园互助群与公益社群的持续互动帖文，排除只有浏览记录、缺少实际回应的材料，以及参与者之间未形成持续联系的一次性争论。在遵守平台公开访问规则并对可识别个人身份的信息作去标识化处理的前提下，结合Python脚本采集与人工定向检索获取公开正文及互动回复，经去重、清理广告灌水与无意义短评并由人工复核后，选取讨论充分、互动度较高的代表性主题帖及其核心回复，按平台与四类价值观维度均衡配置，建成结构化的大学生弱连接交往文本数据库。随后，借助Python开展分词、词频统计、关键词重要性测算与词语共现分析，并严格区分用于检索的场景词（如搭子、打卡群、绩点、上岸）与文本中自然浮现的价值态度词（如消费心态、功利焦虑、边界防御、亲属疏离与公共态度），防止因检索词反复出现而产生同义循环；同时运用短文本主题聚类方法归类不同社区的潜在讨论主题，从宏观上呈现弱连接交往场景与四类价值观表达之间的关联结构。

考虑到大学生在网络社区的发言常常带有反讽、调侃和圈层隐语，单靠计算机自动统计容易误读真实态度，研究在宏观主题归类之后进一步引入质性内容分析法，推进微观话语深描与机理剖析。具体而言，按平台与主题分布分层抽取典型主题帖及其核心回复作为深度分析样本，导入质性分析软件NVivo，依据前期确立的四维价值观界定标准编制操作化编码手册，开展人工精细编码与话语分析，并结合多轮复核校准与跨平台比对控制主观偏倚和单一社区偏差。通过将人工深度编码还原的微观比较心理、交往规则延伸、浅层满足替代与公共责任淡化线索，同宏观统计结果相互印证，系统呈现四类价值观负面偏移的具体表现，拆解四条微观传导路径并剖析背后的经济社会、家庭与高校育人成因；分析过程中明确网络文本资料所能支持的解释范围，立足话语线索辨析与作用过程阐释，不作超出数据支撑的统计因果推论。紧扣前述理论框架、经验证据与成因剖析结果，研究最后通过归纳综合与理论转化，对接新时代高校思想政治教育规律与政策要求，围绕校园实体共同体建设、弱连接资源反向利用、思政课教学转化与校社协同联动，构建兼具针对性与可操作性的高校思想政治教育应对策略。

% 【此处插入技术路线图（如 \begin{figure}[htbp] ... \end{figure}）】

综合来看，上述研究方法与推进路线在理论支撑、语料获取及工具操作上均具备可靠的可行条件。在理论与文献方面，马克思主义经典著作关于社会关系与思想观念的论述、党和国家关于培育社会主义核心价值观和高校思想政治工作的政策文件，以及国内外社会学、传播学与思想政治教育学围绕弱关系、青年轻社交和价值观演变积累的丰硕成果，能够为界定弱连接边界与搭建四维评判标准提供充分依据。在数据来源方面，小红书、豆瓣高校小组、知乎高等教育板块及高校公开论坛汇集了大量在校大学生围绕搭子社交、求职备考互助、亲属往来和校园公共事务的真实记录，通过高校身份核验、双向互动情境限定与分层抽样，能够合规稳定地建成代表性文本数据库，满足宏观统计与人工深度编码的样本要求。在工具与操作方面，Python文本清洗与主题聚类以及NVivo人工内容编码均为社会科学领域成熟可靠的技术手段，各环节衔接紧密、边界清晰，加之前期已具备相关文献积累与文本数据处理基础，能够保障后续实证分析与对策研究顺利落地。
\subsection{研究的创新点与难点}

\subsubsection{研究的创新点}

第一，在研究视角上实现跨学科融通。既有研究多各自展开：社会学与传播学偏重考察弱连接在获取信息、促进求职或缓解焦虑上的实用功能，较少追问其对价值判断尺度的改变；思想政治教育学多侧重主流价值观的宏观阐释和网络不良思潮的内容剖析，较少剖析日常交往方式本身的影响。本研究从马克思在《关于费尔巴哈的提纲》中对人的社会关系的论述出发，把弱连接这一日常交往形式作为切入点，系统剖析其微观作用过程如何冲击马克思主义价值观与社会主义核心价值观的培育，拓宽了思想政治教育学考察青年日常社会关系的研究视角。

第二，在研究内容上细化边界并打通关联场景。本研究以“对象可辨认、有实质回应、有一定延续性”划定弱连接交往的范围，将其与缺少实际回应或持续联系的网络活动区分开来；依据公民道德生活领域确立边界清晰的四维价值观指标，从片面比较、低义务规则延伸、浅层满足替代与松散群体责任淡化四条微观路径揭示作用过程。尤其在人际交往维度，打破了将校园轻社交与家庭关系割裂考察的局限，将校园搭子社交与大学生疏离远亲的现象置于同一交往规则延伸的逻辑中加以阐释，并提出尊重个人边界与重构实体集体相结合的应对思路。

第三，在研究方法上加强微观经验支撑。针对既有思想政治教育研究对日常流动社交话语挖掘相对不足的局限，建立高校身份核验与弱连接交往情境双重限定的多平台网络文本采集流程，将场景词与价值词相分离的宏观统计分析同经过一致性检验的NVivo人工深度编码结合起来，使微观作用过程剖析获得贴近青年真实生活的经验依据。

\subsubsection{研究的难点}

一是抽象价值规范向具体文本编码标准的转化与语境辨识。马克思主义价值观和社会主义核心价值观内涵宏大，而大学生在网络社区中的日常发言零散含蓄，还常夹杂反讽与调侃，需要制定边界清晰、操作性强且便于稳定判别的四维编码细则，确保人工编码的准确性与一致性。

二是从公开网络社区中筛选符合定义的交往材料，并明确这些材料可以支持哪些判断。公开网络社区的用户身份往往难以确认，部分内容也只反映单向浏览或一次性争论，需要在采集阶段同步核验高校在校生身份并严格限定双向互动情境，同时在阐释中明确横截面网络文本能支持的解释范围，避免作出超出经验材料支撑的推论。

% !TeX root = ../main.tex
\section{学位论文工作计划}
\label{sec:work-plan}

\subsection{研究计划时间安排}
\label{sec:schedule}

\begingroup
\small
\renewcommand{\arraystretch}{1.2}
\setlength{\tabcolsep}{5pt}
\renewcommand{\tabularxcolumn}[1]{m{#1}}
\refstepcounter{table}\label{tab:research-schedule}
\noindent
\begin{tabularx}{\textwidth}{|>{\centering\arraybackslash}m{0.06\textwidth}|>{\centering\arraybackslash}m{0.25\textwidth}|>{\centering\setlength{\parindent}{0pt}\arraybackslash}X|}
\hline
序号 & 时间 & 进展计划 \\ \hline
1 & 2025.09-2026.08 & （1）学习专业理论，广泛阅读中外文献\par （2）确定论文选题 \\ \hline
2 & 2026.09-2026.10 & （1）与导师交流研究选题，搭建研究框架\par （2）参加凝练研究方向答辩 \\ \hline
3 & 2026.11-2026.12 & （1）根据凝练研究方向答辩意见进行调整，完善研究框架\par （2）系统梳理国内外文献，完成文献综述\par （3）撰写开题报告，完成开题答辩 \\ \hline
4 & 2027.01-2027.04 & （1）开展资料收集工作，包括发放问卷、开展访谈和进行案例分析等\par （2）将收集到的资料进行整理归类，做好分析，形成相应的表格和数据支撑 \\ \hline
5 & 2027.05-2027.10 & （1）分章节写作，完成论文主体部分\par （2）调整论文格式，打磨文章细节 \\ \hline
6 & 2027.11-2027.12 & （1）根据导师建议修改完善论文\par （2）参加论文中期答辩 \\ \hline
7 & 2028.01-2028.04 & （1）继续撰写学位论文，并根据中期检查意见修改\par （2）进行案例更新、统稿等后期工作\par （3）参加学位论文答辩 \\ \hline
8 & 2028.05-2028.06 & （1）根据专家意见进行修改完善\par （2）完成学位论文终稿，提交归档 \\ \hline
\end{tabularx}
\endgroup

\subsection{已完成的部分工作}
\label{sec:preliminary-work}

\begin{enumerate}[label=\arabic*.,leftmargin=2em,labelsep=0pt]
\item 如期完成“研究计划时间安排”规划的任务1、任务2；
\item 系统梳理国内外关于弱连接、青年价值观及思想政治教育的相关文献，归纳现有研究成果与研究局限，进一步聚焦本研究的核心问题，厘清研究的切入点与创新空间；
\item 已搭建论文整体框架，梳理了弱连接、大学生价值观等核心概念及其理论基础。
\end{enumerate}
